\documentclass[../../main.tex]{subfiles}
\begin{document}

{\bf [解]}
$|a\omega+b|\ge 0$だから題意の式を両辺二乗した
\begin{align}
 |a\omega+b|^2 = 1 \label{eq:1}
\end{align}
を満たす整数$(a,b)$を求めれば良い．
左辺を展開すると
\begin{align}
 |a\omega+b|^2 = a^2|\omega|^2+b^2+ab(\omega+\bar\omega) = a^2+b^2-ab
\end{align}
だから\cref{eq:1}に代入して
\begin{align}
 a^2+b^2-ab=1 \label{eq:2} 
\end{align}
をみたす $(a,b)\in\mathbb{Z}$ をもとめれば良い．
対称性から $a\le b$ とすると\cref{eq:2}より
\begin{align}
a^2+b^2=ab+1\le b^2+1 \quad\therefore\ a^2\le1
\end{align}
だから $a=0,\pm1$ が必要である．
これを順に\cref{eq:2}に代入して，
\begin{align}
(a,b)=(0,\pm1),\ (1,1),\ (1,0),\ (-1,-1),\ (-1,0)
\end{align}
が求める答えである．


\newpage

\end{document}
